% !TeX root = ../../Thesis.tex
\ThesisChapter{\ThesisText{Background}{Grundlagen}}{ch:background}
  {TODO: Add a brief chapter overview. \\
  This chapter explains the domain-specific and technical foundations needed to understand the thesis.}

% TODO: Introduce the technical and conceptual background required to understand the thesis.
%
Explain the fundamental concepts, terminology, technologies, methods, and theoretical principles on which the thesis is based.
The chapter should provide readers with the necessary context to understand the problem, the proposed approach, and the subsequent evaluation.

Possible contents include:
\begin{itemize}
  \item definitions of relevant terms and concepts,
  \item theoretical foundations and established models,
  \item technologies, frameworks, protocols, or algorithms used in the thesis,
  \item the domain-specific context of the research problem,
  \item assumptions and technical constraints that are important for the remainder of the thesis, and
  \item a description of the benchmarks used.
\end{itemize}

Focus only on background information that is directly relevant to the thesis.
Detailed comparisons with existing research and related approaches should normally be presented in the related work chapter.

% Example for numbered equations and for the acronyms declared in src/acronyms.tex.
\section{\ThesisText{Example Definitions}{Beispiel: Definitionen}}
\label{sec:example-definitions}

The first use of an acronym expands it automatically, so \ac{hpc} appears in full here and as a short form from the next use onwards.
Declare each acronym once in \texttt{src/acronyms.tex} and never spell one out by hand.

The speedup of a parallel program on $p$ processing elements follows from the sequential runtime $T(1)$ and the parallel runtime $T(p)$:

\begin{equation}
  \label{eq:speedup}
  S(p) = \frac{T(1)}{T(p)}
\end{equation}

Parallel efficiency relates that speedup to the number of processing elements used to reach it:

\begin{equation}
  \label{eq:efficiency}
  E(p) = \frac{S(p)}{p}
\end{equation}

Equation~\eqref{eq:speedup} and Equation~\eqref{eq:efficiency} are applied to measured data in Chapter~\ref{ch:evaluation}.
A program distributed across several \ac{hpc} nodes with \ac{mpi} is assessed in exactly the same way.
